\chapter{Methodology}

\section{Chapter Introduction}

The purpose of this chapter is to more concretely depict the stages of the pipeline detail what is involved, step by step.
My usage of the word ``pipeline'' is inspired by the nature of the implementation being a full, end-to-end solution (as opposed to specific instances or model developments).
A crucial part of what makes such a pipeline successful is a reliable and operational \emph{flow} of computational tasks.

In a brief sense, we can suggest that there are three main stages to this pipeline, all of which are in terms of underlying source(s) of data:
\begin{enumerate}
    \item \textbf{Acquisition:} finding appropriate data, and preparing it for ``processing'' readiness.
    \item \textbf{Processing:} Otherwise called ``representation'', this step involves transformations of the acquired data to become analytically useful. In ML, for example, this is the training of a model (using the acquired data) to produce meaningful outputs (predictions/classification) that may thereafter be used for analytical purposes.
    \item \textbf{Analytics:} Investigation of underlying data, as represented by the intermediate form, with regards to the inital research question.
\end{enumerate}

Whilst vague, this can help guide expectations for the various steps in between, since with long and complex pipelines it can be difficult to assert the full relevance of each stage in isolation.
A version of this pipeline, therefore, is represented neatly by the diagram in Figure~\ref{fig:pipeline}.

\begin{figure}[ht]
    \hskip-2cm
    \begin{tikzpicture}[
            >=Stealth,
            /tikz/colsep/.initial=4cm,
            /tikz/vertsep/.initial=15mm,
            ghNode/.style={
                    rectangle, draw=black!70, rounded corners,
                    fill=githubPurple!70, anchor = south, text width=3cm, minimum height = 8mm, align=center, font=\small
                },
            jobNode/.style={
                    rectangle, draw=black!70, rounded corners,
                    fill=jobTeal!50, anchor = south, text width=3cm, minimum height = 8mm, align=center, font=\small
                },
            auxNode/.style={
                    rectangle, draw=black!70, rounded corners,
                    fill=auxGray!50, anchor = south, text width=3cm, minimum height = 8mm, align=center, font=\small
                },
            modelNode/.style={
                    rectangle, draw=black!70, rounded corners,
                    fill=modelOrange!50, text width=3cm, minimum height=10mm, align=center, font=\small
                },
            processNode/.style={
                    rectangle, draw=black!70, rounded corners,
                    fill=auxGray!50, text width=3cm, minimum height=10mm, align=center, font=\small
                },
            outcomeNode/.style={
                    rectangle, draw=black!70, rounded corners,
                    fill=outcomeGreen!50, text width=3cm, minimum height=10mm, align=center, font=\small
                },
            evalNode/.style={
                    rectangle, draw=black!70, rounded corners,
                    fill=evalYellow!50, text width=3cm, minimum height=10mm, align=center, font=\small
                },
            % arrow/.style={-Stealth, thick, black!70},
            textArrow/.style={-Stealth, thick, black!70},
            textArrowjobTeal/.style={-Stealth, thick, jobTeal!90},
            textArrowPurple/.style={-Stealth, thick, githubPurple!90},
            quantArrow/.style={-Stealth, thick, black!70, dashed},
            quantArrowjobTeal/.style={-Stealth, thick, jobTeal!90, dashed},
            quantArrowPurple/.style={-Stealth, thick, githubPurple!90, dashed},
            quantArrowEval/.style={-Stealth, thick, evalYellow, dashed},
            % pics/legend entry/.style={code={%   
            %                 \draw[pic actions]
            %                 (-0.5,0.25) sin (-0.25,0.4) cos (0,0.25) sin (0.25,0.1) cos (0.5,0.25);}}
        ]

        \pgfkeysgetvalue{/tikz/colsep}{\colsep}

        \node[auxNode] (aux) at (-\colsep,0) {Auxiliary};
        \node[jobNode] (job)  at (0,0) {Job Postings};
        \node[ghNode] (github) at (\colsep,0) {GitHub Repositories};
        \begin{pgfonlayer}{background}
            \node[draw=black!20, dashed, rounded corners, thick, fill=auxGray!10, fit=(aux) (job) (github), inner sep=4mm]{};
            \node[anchor=south west, font=\small, fill=white] at ([xshift=-15mm, yshift=10mm]aux.north west) {Data Sources};
        \end{pgfonlayer}
        \pgfkeysgetvalue{/tikz/vertsep}{\vsep}
        \node[processNode, below=\vsep of $(github)!0.5!(job)$] (preproc) {Preprocessing};
        \node[modelNode, below=\vsep of preproc] (embed) {Representation};
        \node[modelNode] (weak) at (aux.center |- embed.center) {Weak Baseline};

        \node[modelNode, below=\vsep of embed] (sim) {Similarity};
        \node[evalNode, below=\vsep of weak] (eval) {Evaluation};
        \node[outcomeNode, below=\vsep of sim] (ts) {Time-series};

        \begin{pgfonlayer}{background}
            \node[draw=black!20, dashed, rounded corners, thick, fill=auxGray!10, fit=(weak)(sim) (eval) (embed), inner sep=4mm]{};
            \node[anchor=south west, font=\small, fill=white] at ([xshift=-15mm, yshift=5mm]weak.north west) {NLP Stages};
        \end{pgfonlayer}
        \draw[textArrowjobTeal] (job.south) -- ++(0,-\vsep/8) -| ([xshift=-2mm]preproc.north);
        \draw[textArrowPurple] (github.south) -- ++(0,-\vsep/8) -| ([xshift=2mm]preproc.north);

        \draw[quantArrowEval]
        ([yshift=0.5mm]eval.north east)
        -- ([xshift=-0.5mm]embed.south west)
        -- ([xshift=-0.5mm]sim.north west)
        -- ([xshift=0.5mm]eval.north east);

        % .. controls ([xshift=-3mm, yshift=3mm]eval.north east) and ([xshift=3mm, yshift=3mm]embed.south west) .. (embed.south west)
        % .. controls ([xshift=-3mm, yshift=3mm]embed.south west) and ([xshift=3mm, yshift=-3mm]sim.north west) .. (sim.north west)
        % .. controls ([xshift=-3mm, yshift=3mm]sim.north west) and ([xshift=-3mm, yshift=3mm]eval.north east) .. (eval.north east);

        \draw[textArrow] (aux.south) -- (weak.north);
        % \draw[arrow] (aux.south) -- ++(0,-5mm) |- ++(\colsep/2,0) -- (weak.north);
        \path[textArrowPurple] (github.north) edge [out=90,in=90,looseness=0.5] (aux.north);

        \draw[textArrowjobTeal] ([xshift=-2mm]preproc.south) -- ([xshift=-2mm]embed.north);
        \draw[textArrowPurple] ([xshift=2mm]preproc.south) -- ([xshift=2mm]embed.north);

        \draw[quantArrowjobTeal] ([xshift=-2mm]embed.south) -- ([xshift=-2mm]sim.north);
        \draw[quantArrowPurple] ([xshift=2mm]embed.south) -- ([xshift=2mm]sim.north);

        \draw[quantArrow] ([xshift=-2mm]sim.south) -- ([xshift=-2mm]ts.north);
        \draw[quantArrowPurple]
        ([xshift=2mm]github.south)
        -- ++(0,-\vsep*5.25)
        node[pos = 0.5, right=0.3cm, rotate = 90]{Stars}
        |- ++(-\colsep/2,0)
        -- ([xshift=2mm]ts.north);

        \draw[quantArrow] (weak.south) -- (eval.north);
        \draw[quantArrow] (sim.west) --  (eval.east);
        \matrix [draw, right=0.5*\colsep of embed, font=\tiny, column sep=5mm, row sep=2mm] {
            \draw[-Stealth, thick, black!70] (-0.5,0) -- (0.5,0);                           & \node {\textbf{Textual} data flow}; \\
            \draw[-Stealth, thick, black!70, dashed] (-0.5,0) -- (0.5,0);                   & \node {\textbf{Quantitative} data flow}; \\[3mm]
            \node[circle, draw=jobTeal!70, fill=jobTeal!70, minimum size=4mm] {};           & \node {\textbf{Job-Postings} related};   \\
            \node[circle, draw=githubPurple!70, fill=githubPurple!70, minimum size=4mm] {}; & \node {\textbf{GitHub} related};   \\
            \node[circle, draw=modelOrange!70, fill=modelOrange!70, minimum size=4mm] {};   & \node {\textbf{Model/representation}};           \\
            \node[circle, draw=evalYellow!70, fill=evalYellow!70, minimum size=4mm] {};     & \node {\textbf{Evaluation}};           \\
            \node[circle, draw=outcomeGreen!70, fill=outcomeGreen!70, minimum size=4mm] {}; & \node {\textbf{Outcome}};           \\
            \node[circle, draw=auxGray!70, fill=auxGray!70, minimum size=4mm] {};           & \node {\textbf{Other Process}};           \\
        };

    \end{tikzpicture}
    \caption{Conceptual diagram of methodology pipeline.}
    \label{fig:pipeline}
\end{figure}

Specifically then, if the goal is to measure diffusion of OSS innovation into labor demand, then We want to measure the relationship between two temporal series:
\begin{enumerate}
    \item The trend of \emph{innovation} - The evolution \& popularity on the development side,
    \item and the trend of \emph{adoption} - The requirement of possessing the relevant knowledge/skills.
\end{enumerate}

The first of these is straightforwardly representable by the so-called supply side of development (specifically, \gh{}) metrics.
The second, however, is challenging: there is no straightforward way to depict whether a \repo{} is relevant to a job and vice versa.

Therefore, along the way, we must construct a method to represent this relevance factor.
As the literature suggests, \textbf{Sentence Embeddings} present a clean solution - in an intra-domain sense, the method is demonstrably viable for evaluating the semantic similarity of documents.
In addition, there is evidence to suggest that \emph{cross-domain} semantic similarity can be assessed using sentence embeddings, though there is no research currently that supports this idea in this context (\jobposts{} vs. \ghrepo{}).

The bulk of this paper is therefore dedicated to establishing a \textbf{high quality representation of cross-domain semantic similarity}.
The following describes each component of the pipeline in detail:
\begin{enumerate}
    \item Data collection,
    \item ETL \& EDA (Pre-processing, Exploring \& Cleaning),
    \item Weak model development (Fuzzy-matching keywords),
    \item Sentence embedding representation,
    \item Evaluations and experimentations (Textual Variations, Sentence-transformer variation, Dimensionality Reduction),
    \item and Alignment \& Analysis (Correlation, lead-lag) of resulting time-series'.
\end{enumerate}


\section{Data Collection}

This phase of the project establishes a systemized workflow to find and pull data. This effort can be grouped in three distinct ways categories (or channels):
\begin{enumerate}
    \item \jobposts{},
    \item \ghrepo{},
    \item and Auxiliary data.
\end{enumerate}

\subsection{\jobposts{}}

\subsubsection{Overview of Data Landscape}
As previously discussed, the practice of online advertisment of employment vacancies is a ubiquitous in recent times.
In one sense, this is a benefit because concerns about existence of data are not pertinent.
A panoply of websites, colloquially known as ``job boards'', exist online - the primary examples of these include \emph{LinkedIn}~\cite{linkedin_website} and \emph{Indeed}~\cite{indeed_website}, though these sites are much more than vacancy-advert sites.

Some job boards have dedicated API layers.
For example, among the various available Linkedin API endpoints (under paid access) is one that allows for efficient programmatic interfacing to find vacancies in occupational-fields or specific roles - this is useful for both employers and employees alike.
In other cases, web-scraping is a viable option, albeit more laborious than using APIs.
\subsubsection{Problems in Acquistion}
However, a severe limitation of acquiring data in real-time from these sites is a result of the fleeting nature of \jobposts{}: Once a vacancy is filled, the advert is no longer necessary, so the data-point is no longer maintained on the job board site.
Because of this fact, directly accessing \jobpost{} data from job boards means retrieving that data as it appears (and before it disappears) in real time.

Given that the scope of this paper investigates the long-term relationship between data-sources (diffusion can be a slow process), it is not enough to \emph{begin} collecting data now.

Rather, it is necessary to acquire pre-existing datasets consisting of historic records. This proved problematic, because largely, such data is expensive to store, and as such is not considered by job-board sites as necessary to maintain.
Moreover, because of this economic cost of storing this data, the most comprehensive access to such data is expensive: Firms such as Lightcast~\cite{lightcast__LightcastGlobal} are framed as ``labour market intelligence'' firms that offer data \& analytics services, pertaining to \jobposts{} under the labor market data.

\subsubsection{Solution: Kaggle}
Thankfully, the open data community of Kaggle~\cite{kaggle_datasets} proves mightily useful for such data needs.
This is a platform for Data Science and Machine Learning where people share lessons, data \& code for education and competitive purposes.
As a resource, it is often the solution to data acquisition: in most cases of niche requirements, someone has likely created a relevant dataset on Kaggle.

It turns out that this is true for \jobposts{}, and emphatically so - there exist hundreds of such datasets.
Not all data sets are of equal value - many are very small and many have limited features - And as such I located and downloaded $5$ adequately expressive datasets.
The details of these datasets are expressed in Table~\ref{table:job-data-overview}.

\begin{table}[ht]
    \hskip-2.0cm
    \caption{Overview of datasets used for Job Postings}
    \label{table:job-data-overview}
    \small
    \centering
    \renewcommand{\arraystretch}{1.5}
    \begin{tabular}{p{3.5cm} p{3cm} p{6.5cm} p{2cm}}
        \textbf{Dataset Name}                      & \textbf{Date Range}     & \textbf{Details}                                                  & \textbf{Source}           \\
        \hline
        1000 ML Jobs US                            & 12/2023 - 04/2025       & Mostly within March \& April of 2025, all jobs are ML related     & \cite{__MachineLearninga} \\
        Linkedin Data Engineer Job Postings        & 17/12/2023 - 17/12/2023 & Extra "skills" column - This is appended to "description"         & \cite{asaniczka_2023}     \\
        Data Science Job Postings \& Skills (2024) & 19/01/2024 - 21/01/2024 & Extra "skills" column - This is appended to "description"         & \cite{asaniczka_2024}     \\
        LinkedIn Job Postings (2023 - 2024)        & 23/03/2024 - 20/04/2024 & --                                                                & \cite{arsh_koneru_2024}   \\
        Job Postings from Ireland (October 2022)   & 01/10/2022 - 31/10/2022 & Ireland only; Sample dataset from Techmap.io's commercial product & \cite{__JobPostings}      \\
        US Job Postings from 2023-05-05            & 05/05/2023 - 05/05/2023 & USA Only; Sample dataset from Techmap.io's commercial product     & \cite{__USJoba}           \\
    \end{tabular}
\end{table}

As can be seen, some datasets are created by individuals through web-scraping methods, while others are sample datasets from the much larger stores of commercial entities (e.g. Techmap\cite{techmap}).
Each of these data sets satisfy the minimal criteria for utility in this project. Ranging from a few hundred to a few thousand records, they all possess the required features:
\begin{itemize}
    \item Date component: A timestamp of when the posting was created.
    \item Title component: The advertised job title.
    \item Description component: The main body of texts for a \jobpost{}.
\end{itemize}

\subsubsection{Cleaning \& Collating}

These five datasets cover different temporal (See Fig~\ref{fig:job-data-time-variation}), regional and occupational (See Fig~\ref{fig:job-data-title-variations}) ranges.
This next stage involved reconfiguring their composition into a single representative dataset.
\fbox{\parbox{0.8\linewidth}{\jobposts{}: Distribution of Timestamps}}\label{fig:job-data-time-variation}
\fbox{\parbox{0.8\linewidth}{\jobposts{}: Distribution of Vacancy Titles}}\label{fig:job-data-title-variation}
As is the case for the majority of this project, this process was undertaken using \texttt{Python} (and the library \texttt{Pandas} - full information on utilized software \& libraries is under Appendix~\ref{appendix:software-libaries}).
Some datasets included an additional ``Skills'' field.
During pre-processing, I concatenated this to the ``Description'' field so that all \jobposts{} are represented as a single textual body.
This ensures consistent treatment across postings, whether or not they explicitly separate skills from the main description.

Thereafter, each dataset required varying degrees of processing to retrieve \& represent the $3$ necessary fields, which are thereafter compiled into one big \texttt{Pandas DataFrame}.
Thereafter, I generated a new field - ``Mapped Title'' - which uses \emph{regular expressions} to map the noisy ``Title'' field into a cleaner aggregation field.
Using this, undesirable fields (i.e. not relating to AI or Data Science) were removed - this was done in alignment to the design choice as set out in Section~\ref{section:assumptions}, and the method is largely effective (assuming the proper \texttt{regex} definitions).

As a result, the full \jobposts{} dataset for downstream usage totals $18690$ records ranging from $10/2022 \rightarrow 04/2025$.

\subsection{OSS \repos{}}

\subsubsection{Overview of Data Landscape}

The process for retrieving \ghrepos{} data is much simpler in the availability sense - free API access is offered, via $2$ alternative endpoints: REST API and GraphQL API.
To reiterate, the features of this data channel required are the \repo{} metadata (including the README, description, owner-name, \repo{}-name, and topic-tags) and popularity data, representable by count of stargazers (a `stargazer' of a \repo{} is a user that has starred that \repo{}) as a timeline of the respective creation-events (When the user `starred' the \repo{}).
It is important to note that these are but a subset of the available data points that are collectible, and usable, from \gh{}.

For example, it is possible that other forms of textual data exist on an ad-hoc basis, potentially providing more fine-grained technical information that is relevant to representing its ideas. Furthermore, other measures of popularity and of development evolution Would provide a more holistic understanding of a \repo{}'s true development history.
Nonetheless, the choices made are pragmatic in that they are, on the whole, sufficiently indicative of \repo{} characteristics, for any repository.

\subsubsection{Repository selection}

A subset of \repos{} was required, as acquiring and processing all of \gh{} is infeasible.
Manual curation would be both laborious and subject to bias, so automated alternatives were sought.

I approach this in two different ways:
\begin{enumerate}
    \item Web Scraping by \gh{} topic
    \item Informed by a Third-Party source
\end{enumerate}
The first approach used seed topics (e.g. Machine Learning'', Natural Language Processing''), retrieving the top 30 repositories per topic and recursively expanding via their associated tags.
This method effectively captures a broad subspace of related repositories, but it has drawbacks: niche or irrelevant topics may be included, and many relevant repositories lack topic tags altogether (this latter problem is exactly that as tackled with~\cite{izadi_2021_TopicRecommendation}).

The second approach therefore seeds from an auxiliary data source a list of \repos{} that are trending/relevant to sub-topics of software engineering / AI.
This source, \emph{OSS Insight}~\cite{__OSSInsight}, is an analytics project that provides various functionalities, including daily-updated topical collections of repositories that are \emph{trending} or \emph{popular}.
Their algorithm for determining these repositories include `scores' of stars/forks and a `base score', computing optimized ratings from weighted treatment of such metrics.
Their ranking algorithm is based on a composite score including stars, forks, and a base score, with weights optimized over these metrics. The exact details are not fully disclosed, though their open development process suggests that additional engagement signals may be incorporated (e.g. activity around pull requests and issues, and code evolution in line deletions/creations)~\cite{__WeAre}.
Moreover, their \ghrepo{} reference that `collections' are curated by a \emph{community process}, wherein users can submit collections in pull requests. This constitutes a form of community-driven annotation, providing an effective low-time/effort alternative to developing a bespoke classification mechanism or manually labeling categories.

These \repos{} are externally validated as relevant to AI, reducing the need for subjective judgment on my part. All that is left to me is to select the names of ``Collections'' (analogous to the ``seed topics'' in my first approach) for usage, as seen in table~\ref{table:oss-insight-collections}.
Further consideration is paid to the problems and limitations of this approach are detailed in section~\ref{section:assumptions}.
For the purposes of manageability, this is an appropriate method of sampling a relevant set of \repos{}, and fits nicely into a programmatic pipeline by the simple API access provided.

Importantly, \emph{OSS Insight} metadata (e.g. ‘Collection names’) is not used in downstream analyses, since it references data external to the repository contents. They are perhaps useful for categorization in an evaluative sense, but as they reference data extraneous to the \repo{} contents itself, it would not be appropriate to treat this in kind.

\subsubsection{Acquisition: Web Scraping}
Initially, I chose to develop a simple web scraping method for pulling \repo{} metadata.
Though somewhat effective. This method presents major maintainability issues and furthermore is not necessarily the most compliant mechanism for acquiring website data.

\subsubsection{Acquisition: \gh{}'s REST API}
Thereafter, I graduated to develop a script for acquiring the metadata of \ghrepos{} via \gh{}'s REST API endpoint.
This method is not, however, extensible to observing the star-history for a project. Such data is most effectively stored (and accessed) using a Graph Database, relating as it does to the activity of users, affording the benefit of also storing other forms data inherent to the interconnected network nature of the \gh{} (and broadly, OSS) ecosystem.

\subsubsection{Acquisition: \gh{}'s GraphQL API}
Therefore, I developed a method to access \gh{}'s GraphQL API endpoint, using defined graph-query-language queries that can effectively be managed/altered in a separate part of the codebase.
This method involves paginating through the set of sometimes hundreds of thousands of star-creation events for a \repo{}, extracting the related time-stamps. At $\approx1.5\frac{\text{Seconds}}{\text{Request}}$ \& retrieving $100$ star-events per request, this process is time-intensive for large \repos{} (The largest in this project is \texttt{TensorFlow}'s \repo{}, with $\approx 190,000$ stars).

To consolidate this process, the metadata acquisition script was refactored to also make use the GraphQL API.
This results in a more robust, maintainable and adaptable component for the acquisition of \ghrepo{} data, which remains extensible to acquiring other data-points (e.g. Forks, Commits, Pull Requests, etc.).

\subsubsection{Results}
A full dataset of \ghrepos{} and their related metadata \& stargazer-history are therefore retrieved via maintainable and adaptable code.
The script thereafter aggregates and stores stargazer history in different formats, namely wide and long formats:
\begin{itemize}
    \item \textbf{Long:} A tabular file with fields of \texttt{[`date', `repository', `stars']}.
    \item \textbf{Wide:} A tabular file resembling a matrix of shape $(n, m)$, where $n$ is the number of repositories and $m$ is the number of discrete time periods (daily, monthly, or quarterly).
          The time axis is defined as the range from the earliest star-event across all repositories ($\min D$) to the latest ($\max D$), partitioned into equal-length periods. Each entry $(i, t)$ records the number of star-events for repository $i$ in period $t$.
\end{itemize}

\subsection{Auxiliary Resources}
The term `Auxiliary' here represents a channel of data that is both an assortment of sources that were used for evaluative (or otherwise informative) purposes and not central to the representation/analysis itself.
\subsubsection{Keyword Dictionary}
As is made clear in section~\ref{section:weak-model-specification}, there is a need for some `ground truth' measure against which representation methods can be evaluated.

In essence, there is a requirement for such abstract representation methods for \emph{labelled} data that is not easily devised.
This would entail definitively creating a metric of ``similarity'' between \jobpost{}/\repo{} pairs $(R_r, J_j)$\footnote{To aid notation and simpler discussion, an ambiguous pair of (\jobpost{}/\repo{}) is $(R_r, J_j)$, where $R_r$ and $J_j$ refer to the $r^{\text{th}}$ and $j^{\text{th}}$ \repo{} and \jobpost{} data-points respectively},
a task that is difficult to define clearly, as well as being heavily time-intensive and laborious.
To overcome this issue, a heuristic based method was developed that utilizes a basic \emph{dictionary} of relevant terms/phrases for keyword matching.
The curation of this data is integral to the methodology of the evaluation baseline (hereafter referred to as the \texttt{WeakModel}), so the details are left to section~\ref{section:weak-model-specification}.

\subsubsection{ESCO Skill Classification}
The full ESCO Skills \& Competences classification dataset was downloaded and used for exploring the contents of \jobposts{}.

This data is not utilized in any meaningful sense by the project's pipeline - this is a methodological alignment with philosophy such as in~\cite{escoe-2018-taxonomy},
expressing a preference to rely on a data-driven \& discovered taxonomy. While there is conceivable utility to the ESCO classification,
it is not immediately clear that the sort of shared semantics between \jobposts{} and \ghrepos{} would resemble their taxonomy, so their contents are left aside for this paper's primary developments.

\section{Weak Model Baseline}\label{section:weak-model-specification}

In supervised learning, evaluation is straightforward because we have a \textbf{target} or ground truth: we know the correct labels or outcomes, so we can measure how well our model’s predictions match reality.
For similarity analysis, however, the notion of a target is less obvious. We need to define what it means for two items to be “similar” and establish a reference or standard against which our representation method can be judged.

Intuitively, this means constructing a set of examples (such as pairs of texts) where the degree of similarity is known or agreed upon—either through expert annotation, crowd-sourced judgments, or some other reliable method.
These examples serve as the \textbf{benchmark} or target, allowing for systematic comparison of different representation methods to determine which one best captures the desired notion of similarity.

In this case, the problem lies with labels - there are none. It is therefore necessary to define labels before proceeding to evaluating the representation methods.
The chosen method follows, broadly, fuzzy matching a dictionary of hand-crafted keywords.

\subsection{Keyword Dictionary Design}


\begin{itemize}
    \item Construction of seed keyword list.
    \item Curation and refinement process.
    \item Matching rules.
\end{itemize}

\subsection{Job–Repo Pair Labeling}
\begin{itemize}
    \item Description of how weak model generates positive/negative matches.
\end{itemize}

\subsection{Purpose of Baseline}
\begin{itemize}
    \item Provides pseudo-ground truth for evaluation.
\end{itemize}

\section{Representation Methods}
\subsection{Text Preprocessing}
\begin{itemize}
    \item Tokenization and filtering.
    \item Field selection (job titles vs repo descriptions).
\end{itemize}

\subsection{Sentence Embedding Models}
\begin{itemize}
    \item Model choice (e.g., transformer-based embeddings).
    \item Input fields for embeddings.
\end{itemize}

\subsection{Similarity Computation}
\begin{itemize}
    \item Cosine similarity between job and \repo{} vectors.
    \item Production of job–repo similarity matrix.
\end{itemize}

% \begin{figure}[h]
%     \centering
%     \includegraphics[width=0.7\textwidth]{figures/similarity_matrix.pdf}
%     \caption{Illustration of job–repository similarity matrix construction.}
%     \label{fig:similarity_matrix}
% \end{figure}
\fbox{\parbox{0.8\linewidth}{Illustration of job–repository similarity matrix construction.}}

\section{Time-Series Construction}
\subsection{\repo{} Popularity}
\begin{itemize}
    \item Quarterly star count aggregation.
    \item Handling of older vs newer \repos{} (zero-padding).
\end{itemize}

\subsection{Job–Repo Similarity Series}
\begin{itemize}
    \item Aggregation of similarity values across time.
\end{itemize}

\subsection{Alignment of Series}
\begin{itemize}
    \item Aligning job–repo similarity time series with \repo{} popularity series.
    \item Handling mismatched time horizons.
\end{itemize}

\section{Evaluation Metrics}

\begin{itemize}
    \item Correlation coefficient: degree of linear relationship.
    \item Maximum cross-correlation and lag: lead–lag structure between series.
    \item Fraction of high-alignment periods: consistency of strong similarity.
    \item Popularity-adjusted similarity: weighting similarity by repo star counts.
\end{itemize}

% \begin{figure}[h]
%     \centering
%     \includegraphics[width=0.75\textwidth]{figures/metrics_example.pdf}
%     \caption{Example illustration of correlation and lag metrics between \repo{} popularity and job similarity series.}
%     \label{fig:metrics_example}
% \end{figure}
\fbox{\parbox{0.8\linewidth}{Example illustration of correlation and lag metrics between \repo{} popularity and job similarity series.}}

\section{Summary}
\begin{itemize}
    \item Recap of pipeline.
    \item How each component contributes to testing the research question.
    \item Transition to Results/Analysis.
\end{itemize}
